\begin{corollary}[K17-K32-ONEPRIME]\label{cor:k32-oneprime}
Under the stipulated arithmetic theorem of Alp\"oge,
$K_{32}(\mathbf Z)$ is a finite $q$-primary group, where
\[
                             q=25078104767246023.
\]
Equivalently, its $\ell$-primary subgroup is zero for every prime
$\ell\ne q$. This conclusion allows the group itself to be zero
and gives no numerical bound on its order.
\end{corollary}

\begin{proof}
The group is finite by the established finiteness theorem for the
positive $K$-groups of $\mathbf Z$ in degrees divisible by four.
For an odd prime $\ell<2^{31}$, the published Vandiver verification
and (K2)--(K3) give $K_{32}(\mathbf Z)\{\ell\}=0$.
The 2-primary subgroup is zero by Rognes--Weibel, as recorded after
(K3).

Now let $\ell>2^{31}$ with $\ell\ne q$ and set $e=16$.
The same closed datum
$\mathrm{GD}(\Omega,\mathbf Q(\sqrt{-3}),\ell,16,1)$ at modulus two
is valid: $e$ is positive and even, $\ell\ne2,3$, and
$\ell-1>32$.
The second row of (HH.2) has $q$ as its only prime factor exceeding
the published verification range.
Hence at least one of $N_{16,7}$ and $N_{16,13}$ is an $\ell$-adic
unit. Its auxiliary prime is split and differs from $\ell$.
The uniform full-space coverage and conductor bridge therefore
satisfy Corollary~\ref{cor:general-hecke-unit}, which makes the
strict negative-character Selmer group zero.
Equation (K5) gives $K_{32}(\mathbf Z)\{\ell\}=0$.
There is no prime equal to $2^{31}$, so these cases exhaust all
$\ell\ne q$ and prove the assertion.
\end{proof}
